\scanpage{148}
et que $\int_0^{+\infty}\frac{n(\unsure{r|v})}{\unsure{r|v}^{1-p}}\,d\unsure{r|v}<+\infty$, à fortiori $\int_{\unsure{1|0}}^{+\infty}\frac{n(\unsure{r|v})}{r}\,d\unsure{r|v}<+\infty$,

donc $p\rho^\alpha\int_\rho^{+\infty}\frac{n(\unsure{r|v})}{r}\,d\unsure{r|v}<p\rho^\alpha\int_0^{+\infty}\frac{n(\unsure{r|v})}{r}\,dr\to0$ si $\rho\to0$.

\par\noindent\unsure{\rule{0.95\linewidth}{0.2pt}}\par

\deleted{pour \unsure{tout|tt} $\varepsilon$, on peut \unsure{supposer} \unsure{aussi} que}

\[
\deleted{\sum\|A_{\deleted{\unsure{i}}\unsure{i+1}}-A_i\|_p<+\infty,\quad
\sum S_p(A_{\unsure{i-1|i+1}}-A)<+\infty}
\]

\deleted{alors $A=\sum A_{i+1}-A_i$, \unsure{d'où|donc} $A\leq\sum|A_{i+1}-A_i|$}

\uline{Preuve de b} On a

\[
\begin{aligned}
\sum_{r_\nu(A_{\unsure{i|n}})\leq\rho}(\unsure{r|v}_\nu(A_{\unsure{i|n}}))^{\unsure{r|v|\nu|p}}&=S_p(A_{\unsure{i|n}})-\sum_{r_\nu(A_{\unsure{i|n}})>\rho}(r_\nu(A_{\unsure{i|n}}))^p\\
&\longrightarrow\deleted{\unsure{\sum|\leq}} S_p(A)-\sum_{r_\nu\unsure{>|\geq}\rho}r_\nu^p=\sum_{r_\nu\unsure{\leq|<}\rho}r_\nu^p
\end{aligned}
\]

il suffit de prendre $\rho$ tel que $\sum_{r_\nu\leq\rho}\unsure{r|v}_\nu^{\unsure{p|r}}<\varepsilon$, puis $i$ assez grand.
